\documentclass[11pt,a4paper]{report}

% ---- Packages (all common; compiles with pdflatex) ----------------------
\usepackage[utf8]{inputenc}
\usepackage[T1]{fontenc}
\usepackage[margin=1in]{geometry}
\usepackage{parskip}
\usepackage{xcolor}
\usepackage{listings}
\usepackage{enumitem}
\usepackage{longtable}
\usepackage{titlesec}
\usepackage[hidelinks]{hyperref}

% ---- Listings style ------------------------------------------------------
\definecolor{codebg}{RGB}{245,246,249}
\definecolor{codekw}{RGB}{37,99,235}
\definecolor{codecomment}{RGB}{107,114,128}
\definecolor{codestr}{RGB}{21,128,61}

\lstset{
  basicstyle=\ttfamily\footnotesize,
  backgroundcolor=\color{codebg},
  keywordstyle=\color{codekw}\bfseries,
  commentstyle=\color{codecomment}\itshape,
  stringstyle=\color{codestr},
  showstringspaces=false,
  breaklines=true,
  columns=fullflexible,
  frame=single,
  framesep=4pt,
  rulecolor=\color{gray!30},
  numbers=none,
  tabsize=2,
  aboveskip=8pt,
  belowskip=8pt,
}

\newcommand{\qa}[3]{%
  \textbf{Ideal.} #1\par
  \textbf{Common mistakes.} #2\par
  \textbf{Follow-up.} #3\par
  \vspace{4pt}
}

\titleformat{\chapter}[display]{\normalfont\huge\bfseries}{\chaptertitlename\ \thechapter}{20pt}{\Huge}

% ---- Title ---------------------------------------------------------------
\title{\textbf{Role-Based Access Control (RBAC) Service}\\[6pt]
\large Technical Report and Interview Preparation Guide}
\author{Take-Home Assignment}
\date{}

\begin{document}

\maketitle

\begin{center}
\textit{This report is not for grading. It is a study companion so the author can
confidently explain, defend, and extend every part of the RBAC service during a
technical interview.}
\end{center}

\tableofcontents

% =========================================================================
\chapter{Introduction}
% =========================================================================

\section{What RBAC Is}

Role-Based Access Control (RBAC) is an authorization model that decides
\emph{who can do what}. Rather than attaching permissions to individual users,
RBAC introduces an intermediate concept---the \emph{role}. Permissions are
granted to roles, and roles are granted to users.

\begin{verbatim}
   User  --- has --->  Role  --- has --->  Permission (resource + action)
\end{verbatim}

There are three core entities:

\begin{itemize}[leftmargin=1.4em]
  \item \textbf{Permission}: the atomic unit of access, a pair of
        \emph{resource} (what, e.g. ``documents'') and \emph{action}
        (how, e.g. ``read'' or ``write'').
  \item \textbf{Role}: a named bundle of permissions (e.g. ``admin'').
  \item \textbf{User}: a person or system account that is granted one or more
        roles.
\end{itemize}

The central operation is a boolean check:
\texttt{can\_user\_perform\_action(user\_id, resource, action)}.

\section{Why Companies Use RBAC}

\begin{itemize}[leftmargin=1.4em]
  \item \textbf{Simplicity at scale.} Managing permissions per user does not
        scale. With roles, you assign a handful of roles instead of dozens of
        individual permissions.
  \item \textbf{Consistency.} Everyone with the same role has exactly the same
        access. There is no drift between users who ``should'' be equivalent.
  \item \textbf{Auditability.} ``Who can delete invoices?'' becomes ``which
        roles include invoices:delete, and who has those roles?''---an easy
        query.
  \item \textbf{Least privilege.} Roles encode job functions, making it natural
        to give people only the access their job requires.
  \item \textbf{Change management.} Update a role once and every holder is
        updated instantly. Onboarding/offboarding is assigning/removing roles.
\end{itemize}

\section{Real-World Examples}

\begin{itemize}[leftmargin=1.4em]
  \item \textbf{Cloud platforms} (AWS IAM, GCP IAM): roles/policies grant
        permissions on resources.
  \item \textbf{Operating systems}: groups (roles) grant file permissions.
  \item \textbf{Databases}: \texttt{GRANT}/\texttt{REVOKE} on roles.
  \item \textbf{SaaS products}: Owner / Admin / Member / Viewer tiers.
  \item \textbf{Version control} (GitHub): organisation and repository roles.
\end{itemize}

\section{Scope of This Project}

This project implements a compact but complete RBAC service: an in-memory data
model, a service layer with the permission-checking algorithm, a REST API, seed
data, a demo frontend, and a test suite. Authentication, databases, and
infrastructure are deliberately out of scope and are discussed only in the
scalability and feature-extension chapters.

% =========================================================================
\chapter{Architecture}
% =========================================================================

\section{Layered Design}

The backend is split into layers, each with a single responsibility:

\begin{verbatim}
   Frontend (HTML/CSS/JS)
         |   HTTP + JSON
         v
   routes.py       ->  HTTP concerns: validate input, format output
         v
   services.py     ->  Business logic: RBACService + permission check
         v
   storage.py      ->  Persistence: InMemoryStore (plain dicts)
\end{verbatim}

This separation is the single most important design decision in the project.
The service layer knows nothing about HTTP, and the storage layer knows nothing
about business rules. Each layer can be understood, tested, and replaced on its
own.

\section{Folder and File Responsibilities}

\begin{longtable}{p{0.28\textwidth} p{0.66\textwidth}}
\textbf{Path} & \textbf{Responsibility} \\
\hline
\endhead
\texttt{app/models.py} & Domain objects: \texttt{Permission}, \texttt{Role},
\texttt{User} as dataclasses. Pure data, no logic. \\
\texttt{app/storage.py} & \texttt{InMemoryStore}: two dictionaries keyed by id
(\texttt{users}, \texttt{roles}). \\
\texttt{app/services.py} & \texttt{RBACService}: all create/link/edit operations
and \texttt{can\_user\_perform\_action}. Also \texttt{NotFoundError}. \\
\texttt{app/schemas.py} & Pydantic request/response models---the wire contract. \\
\texttt{app/routes.py} & FastAPI endpoints; thin adapters over the service. \\
\texttt{app/dependencies.py} & Builds the seeded service and exposes
\texttt{get\_service} for \texttt{Depends()}. \\
\texttt{app/seed.py} & Creates example roles, permissions, and users. \\
\texttt{app/main.py} & App creation, CORS, 404 handler, router mount, static
frontend. \\
\texttt{frontend/} & \texttt{index.html}, \texttt{styles.css}, \texttt{script.js}. \\
\texttt{tests/test\_rbac.py} & Service-level and API-level pytest tests. \\
\end{longtable}

\section{Why FastAPI}

\begin{itemize}[leftmargin=1.4em]
  \item \textbf{Validation for free} via Pydantic: request bodies are parsed and
        validated from type hints; invalid input yields a 422 automatically.
  \item \textbf{Automatic docs}: OpenAPI schema and Swagger UI at
        \texttt{/docs}, ReDoc at \texttt{/redoc}.
  \item \textbf{Dependency injection} via \texttt{Depends}, which we use to hand
        the service to routes and to override it in tests.
  \item \textbf{Async-ready and fast}, built on Starlette and Uvicorn (ASGI).
  \item \textbf{Minimal boilerplate}: decorated functions become endpoints.
\end{itemize}

\section{Why In-Memory Storage}

The assignment requires in-memory storage. It has real benefits for this scope:
zero setup, no external dependencies, trivially fast, and it keeps attention on
the RBAC logic rather than on database plumbing.

\textbf{Trade-offs.} Data is lost on restart; it is limited to a single process
(so it cannot scale horizontally without a shared store); and there are no
transactions or durability. These are acceptable for a demo and are exactly the
first things addressed in the scalability chapter.

\section{Design Trade-offs Summary}

\begin{itemize}[leftmargin=1.4em]
  \item \textbf{References vs. embedding}: users store \texttt{role\_ids}, not
        \texttt{Role} objects, so a role is a single source of truth.
  \item \textbf{Sets vs. lists}: permissions and role ids are stored in
        \texttt{set}s for O(1) membership and free de-duplication.
  \item \textbf{Thin storage vs. rich repository}: the store is just two dicts;
        logic lives in the service. Simple to explain; the seam still allows a
        later swap to a database.
\end{itemize}

% =========================================================================
\chapter{Data Model}
% =========================================================================

\section{Entities}

\begin{lstlisting}[language=Python]
@dataclass(frozen=True)
class Permission:
    resource: str   # e.g. "documents"
    action: str     # e.g. "read" or "write"

@dataclass
class Role:
    id: str
    name: str
    permissions: set[Permission]

@dataclass
class User:
    id: str
    name: str
    role_ids: set[str]
\end{lstlisting}

\section{Why These Choices}

\begin{itemize}[leftmargin=1.4em]
  \item \textbf{\texttt{Permission} is \texttt{frozen}}: an immutable dataclass
        is hashable, so it can be an element of a \texttt{set}. Two permissions
        with the same resource and action are equal and de-duplicated.
  \item \textbf{\texttt{Role.permissions} is a \texttt{set}}: adding a duplicate
        permission is a no-op; membership test is O(1).
  \item \textbf{\texttt{User.role\_ids} is a \texttt{set} of ids}: references
        keep one copy of each role; assigning the same role twice is a no-op.
\end{itemize}

\section{Relationships}

\begin{verbatim}
   +--------+          +--------+          +--------------------+
   |  User  | 1 ---- * |  Role  | 1 ---- * |     Permission     |
   +--------+          +--------+          | (resource, action) |
   | id     |          | id     |          +--------------------+
   | name   |          | name   |
   | roles* | -------> | perms* | -------> stored as a set
   +--------+  by id   +--------+
\end{verbatim}

A user has many roles; a role has many permissions. Both relationships are
many-to-many in general (a role can be held by many users; a permission can
appear in many roles), which is exactly why RBAC scales better than
per-user permissions.

\section{Identity}

Ids are server-generated UUID4 strings. Names (``Alice'', ``admin'') are human
labels, not identifiers, so two roles could share a name without colliding. The
API always operates on ids.

% =========================================================================
\chapter{API Design}
% =========================================================================

Base URL: \texttt{http://127.0.0.1:8000}. All bodies are JSON.

\section{Endpoint Reference}

\begin{longtable}{p{0.09\textwidth} p{0.42\textwidth} p{0.34\textwidth}}
\textbf{Method} & \textbf{Path} & \textbf{Purpose} \\
\hline
\endhead
POST & \texttt{/users} & Create a user (optional roles) \\
GET & \texttt{/users} & List users \\
POST & \texttt{/users/\{user\_id\}/roles} & Assign role \\
DELETE & \texttt{/users/\{user\_id\}/roles/\{role\_id\}} & Remove role \\
DELETE & \texttt{/users/\{user\_id\}} & Delete user (204) \\
POST & \texttt{/roles} & Create a role (optional perms) \\
GET & \texttt{/roles} & List roles \\
POST & \texttt{/roles/\{role\_id\}/permissions} & Add permission \\
DELETE & \texttt{/roles/\{role\_id\}/permissions} & Remove permission (body) \\
DELETE & \texttt{/roles/\{role\_id\}} & Delete role, cascades (204) \\
POST & \texttt{/check-permission} & Check a permission \\
\end{longtable}

\section{Per-Endpoint Details}

\subsection*{POST /users}
\textbf{Purpose.} Create a user, optionally with initial roles.\\
\textbf{Request.} \texttt{\{"name": "Alice", "role\_ids": ["..."]\}}\\
\textbf{Response.} \texttt{201} with \texttt{\{id, name, role\_ids\}}.\\
\textbf{Validation.} \texttt{name} required, non-empty; \texttt{role\_ids}
defaults to \texttt{[]}.\\
\textbf{Errors.} \texttt{404} if any \texttt{role\_id} does not exist;
\texttt{422} if \texttt{name} is missing.

\subsection*{GET /users}
\textbf{Purpose.} List all users. \textbf{Response.} \texttt{200} with an array
of user objects. \textbf{Errors.} none.

\subsection*{POST /users/\{user\_id\}/roles}
\textbf{Purpose.} Grant a role to a user.\\
\textbf{Request.} \texttt{\{"role\_id": "..."\}}. \textbf{Response.} \texttt{200}
with the updated user.\\
\textbf{Errors.} \texttt{404} if the user or role does not exist. Idempotent:
re-assigning the same role is a no-op.

\subsection*{DELETE /users/\{user\_id\}/roles/\{role\_id\}}
\textbf{Purpose.} Revoke a role.\\
\textbf{Response.} \texttt{200} with the updated user. \textbf{Errors.}
\texttt{404} if the user does not exist. Removing a role the user does not have
succeeds (idempotent DELETE).

\subsection*{POST /roles}
\textbf{Purpose.} Create a role with optional permissions.\\
\textbf{Request.} \texttt{\{"name": "editor", "permissions": [\{"resource":
"documents", "action": "write"\}]\}}.\\
\textbf{Response.} \texttt{201} with \texttt{\{id, name, permissions\}}.
\textbf{Errors.} \texttt{422} on invalid body.

\subsection*{GET /roles}
\textbf{Purpose.} List all roles with their permissions. \textbf{Response.}
\texttt{200} with an array of role objects.

\subsection*{POST /roles/\{role\_id\}/permissions}
\textbf{Purpose.} Add one permission to a role.\\
\textbf{Request.} \texttt{\{"resource": "documents", "action": "write"\}}.
\textbf{Response.} \texttt{200} with the updated role. \textbf{Errors.}
\texttt{404} if the role does not exist. Idempotent.

\subsection*{DELETE /roles/\{role\_id\}/permissions}
\textbf{Purpose.} Remove one permission from a role. The permission to remove is
sent in the request body (there is no natural id for a permission).\\
\textbf{Request.} \texttt{\{"resource": "documents", "action": "write"\}}.
\textbf{Response.} \texttt{200} with the updated role. \textbf{Errors.}
\texttt{404} if the role does not exist. Idempotent.

\subsection*{POST /check-permission}
\textbf{Purpose.} The core question: may this user do this?\\
\textbf{Request.} \texttt{\{"user\_id": "...", "resource": "...", "action":
"..."\}}. \textbf{Response.} \texttt{200} with \texttt{\{"allowed": true|false\}}.\\
\textbf{Design note.} An unknown user returns \texttt{\{"allowed": false\}}
rather than a 404, matching the specification (``return True if the user has the
permission, False otherwise'').

\section{Status Codes and Error Handling}

\begin{itemize}[leftmargin=1.4em]
  \item \texttt{200 OK}---successful reads and updates.
  \item \texttt{201 Created}---resource created (users, roles).
  \item \texttt{204 No Content}---resource deleted (empty body). Deleting a role
        cascades: it is discarded from every user's \texttt{role\_ids}.
  \item \texttt{404 Not Found}---raised by the service as \texttt{NotFoundError}
        and converted centrally by an exception handler in \texttt{main.py}.
  \item \texttt{422 Unprocessable Entity}---Pydantic validation failure,
        produced automatically by FastAPI.
\end{itemize}

% =========================================================================
\chapter{Complete Code Walkthrough}
% =========================================================================

\section{models.py}

Defines the three domain dataclasses (see the Data Model chapter). No functions
beyond the generated \texttt{\_\_init\_\_}/\texttt{\_\_eq\_\_}/\texttt{\_\_hash\_\_}.
The single subtlety is \texttt{frozen=True} on \texttt{Permission} to make it
hashable and set-storable.

\section{storage.py}

\begin{lstlisting}[language=Python]
class InMemoryStore:
    def __init__(self) -> None:
        self.users: dict[str, User] = {}
        self.roles: dict[str, Role] = {}
\end{lstlisting}

\textbf{Purpose.} Hold state. \textbf{Complexity.} Dictionary lookups by id are
O(1) average. That is all the storage layer does.

\section{services.py}

\subsection*{NotFoundError}
A small exception carrying the entity type and id. Raised by lookups; translated
to HTTP 404 in \texttt{main.py}.

\subsection*{create\_role(name, permissions)}
\textbf{Params.} name, optional list of permissions. \textbf{Returns.} the new
\texttt{Role}. \textbf{Logic.} generate a UUID, wrap permissions in a set, store
by id. \textbf{Complexity.} O(P) to build the set.

\subsection*{get\_role(role\_id) / get\_user(user\_id)}
\textbf{Logic.} dictionary lookup; raise \texttt{NotFoundError} if missing.
\textbf{Complexity.} O(1).

\subsection*{add\_permission / remove\_permission}
\textbf{Logic.} fetch the role (404 if missing), then \texttt{set.add} or
\texttt{set.discard}. Both are idempotent. \textbf{Complexity.} O(1).

\subsection*{create\_user(name, role\_ids)}
\textbf{Logic.} validate that every referenced role exists (calling
\texttt{get\_role}, which raises 404 on the first missing one), then create and
store the user. \textbf{Complexity.} O(R) in the number of roles supplied.

\subsection*{assign\_role / remove\_role}
\textbf{Logic.} fetch the user (404 if missing); \texttt{assign\_role} also
validates the role exists, then \texttt{set.add}; \texttt{remove\_role} uses
\texttt{set.discard}. \textbf{Complexity.} O(1).

\subsection*{can\_user\_perform\_action(user\_id, resource, action)}
The core algorithm; see the next chapter for the full step-by-step and
flowchart. \textbf{Complexity.} O(R) in the number of a user's roles.

\section{schemas.py}

Pydantic models mirror the request/response shapes. \texttt{Field(min\_length=1)}
rejects empty strings; \texttt{default\_factory=list} makes list fields optional
with a safe default. These schemas are the only place input is trusted after
validation.

\section{routes.py}

Each endpoint: (1) receives a validated Pydantic body and/or path params, (2)
receives the service through \texttt{Depends(get\_service)}, (3) calls one
service method, (4) maps the returned domain object to a response schema via the
helpers \texttt{\_user\_to\_response} / \texttt{\_role\_to\_response} (which sort
their output so responses are deterministic). Routes contain no business logic.

\section{dependencies.py}

Builds one seeded \texttt{RBACService} at import time and exposes
\texttt{get\_service()}. Routes depend on this function, so tests can replace it
with \texttt{app.dependency\_overrides[get\_service]}.

\section{seed.py}

Creates the four permissions, the three roles (admin, editor, viewer), and the
three users (Alice, Bob, Charlie), wiring them per the example scenario.

\section{main.py}

Creates the \texttt{FastAPI} app, enables permissive CORS for local development,
registers the \texttt{NotFoundError} to 404 handler, includes the router, and
mounts the \texttt{frontend/} directory as static files at \texttt{/}. API
routes are registered before the static mount so they take precedence.

% =========================================================================
\chapter{Permission Checking Algorithm}
% =========================================================================

\section{The Function}

\begin{lstlisting}[language=Python]
def can_user_perform_action(self, user_id, resource, action) -> bool:
    user = self.store.users.get(user_id)
    if user is None:
        return False
    required = Permission(resource=resource, action=action)
    for role_id in user.role_ids:
        role = self.store.roles.get(role_id)
        if role is not None and required in role.permissions:
            return True
    return False
\end{lstlisting}

\section{Step by Step}

\begin{enumerate}[leftmargin=1.6em]
  \item \textbf{Find the user.} A missing user can do nothing, so return
        \texttt{False} immediately.
  \item \textbf{Build the target permission.} Construct
        \texttt{Permission(resource, action)}. Because \texttt{Permission} is a
        frozen dataclass, equality and hashing are by value.
  \item \textbf{Scan roles.} For each role id on the user, look up the role and
        test whether the target permission is in its permission set.
  \item \textbf{Short-circuit.} The first role that grants it returns
        \texttt{True}.
  \item \textbf{Default deny.} If no role grants it, return \texttt{False}.
\end{enumerate}

\section{Flowchart}

\begin{verbatim}
        +-------------------------+
        |   user_id, resource,    |
        |         action          |
        +-----------+-------------+
                    |
                    v
        +-------------------------+     no
        |  user exists in store?  +----------> return False
        +-----------+-------------+
                    | yes
                    v
        +-------------------------+
        | required = Permission(  |
        |   resource, action)     |
        +-----------+-------------+
                    |
                    v
        +-------------------------+
        | for each role_id of the |
        |          user           |
        +-----------+-------------+
                    |
                    v
        +-------------------------+    yes
        | required in role.perms? +----------> return True
        +-----------+-------------+
                    | no (next role)
                    v
        +-------------------------+
        |  roles exhausted        +----------> return False
        +-------------------------+
\end{verbatim}

\section{Complexity}

Let $R$ be the number of roles a user has and each membership test be O(1) on a
set. The check is \textbf{O($R$)} time and O(1) extra space. In practice $R$ is
small (single digits), so checks are effectively constant time.

\section{Optimizations (Discussion Only)}

\begin{itemize}[leftmargin=1.4em]
  \item \textbf{Precompute an effective permission set per user} (the union of
        all their roles' permissions). Then a check is a single O(1) set lookup.
        Trade-off: must invalidate/recompute when roles or permissions change.
  \item \textbf{Cache in Redis} keyed by user id; invalidate on writes.
  \item \textbf{Bitmasks} for a small, fixed permission universe: represent each
        role's permissions as a bitmask and OR them; a check is a bit test.
  \item \textbf{Index by permission} (permission -> roles -> users) if the common
        query is ``who can do X''.
\end{itemize}

% =========================================================================
\chapter{FastAPI Concepts}
% =========================================================================

\section{Routers}
\texttt{APIRouter} groups related endpoints; \texttt{app.include\_router(router)}
mounts them. This keeps \texttt{main.py} small and lets endpoints live in their
own module.

\section{Pydantic}
Pydantic v2 models declare fields with type hints. On input, FastAPI parses JSON
into the model, coercing and validating types; on output, \texttt{response\_model}
serialises and filters the returned data to the declared schema.

\section{Validation}
Constraints like \texttt{Field(min\_length=1)} are enforced automatically. A
failing request never reaches your function; FastAPI returns a structured
\texttt{422} describing what was wrong.

\section{Dependency Injection}
\texttt{Depends(get\_service)} tells FastAPI to call \texttt{get\_service} and
pass its result as an argument. Dependencies can be nested and overridden. We use
one dependency to supply the service, which makes tests swap in a clean store
with \texttt{app.dependency\_overrides}.

\section{Response Models}
Declaring \texttt{response\_model=UserResponse} (or \texttt{list[UserResponse]})
documents the endpoint in OpenAPI and guarantees the response shape, hiding any
internal fields not in the schema.

\section{HTTP Status Codes}
Set per route with \texttt{status\_code=status.HTTP\_201\_CREATED}, or raised via
exceptions. FastAPI defaults to \texttt{200} for successful responses.

% =========================================================================
\chapter{Project Flow}
% =========================================================================

A single permission check travels through every layer:

\begin{verbatim}
  1. Frontend  : user picks User/Resource/Action, clicks "Check Permission".
                 script.js POSTs JSON to /check-permission.
  2. API       : FastAPI validates the body against PermissionCheckRequest.
                 Invalid -> 422 before any code runs.
  3. Route     : check_permission() receives the body and the service
                 (via Depends) and calls the service method.
  4. Service   : can_user_perform_action() looks up the user, builds the
                 target Permission, scans the user's roles' permission sets.
  5. Storage   : dict lookups by id return the User and Role objects.
  6. Response  : the boolean is wrapped in PermissionCheckResponse and
                 serialised to {"allowed": true|false}.
  7. Frontend  : renders "Allowed" or "Denied".
\end{verbatim}

The same path applies to writes (create/assign/add), except the service mutates
the store and the route returns the updated resource. Errors raised in the
service (\texttt{NotFoundError}) are caught by the exception handler and become
404 responses.

% =========================================================================
\chapter{Testing}
% =========================================================================

\section{Strategy}
Two levels: fast \emph{service-level} unit tests that exercise the logic without
HTTP, and a few \emph{API-level} tests using FastAPI's \texttt{TestClient} with a
fresh service injected through \texttt{dependency\_overrides}. Fixtures build a
clean \texttt{RBACService(InMemoryStore())} per test, so tests are isolated.

\section{Test-by-Test}

\begin{longtable}{p{0.42\textwidth} p{0.54\textwidth}}
\textbf{Test} & \textbf{Why it exists} \\
\hline
\endhead
permission\_check\_allowed & Positive path: a granted permission returns True. \\
permission\_check\_denied & Negative path: a missing permission returns False. \\
assign\_role\_grants\_permission & Assigning a role changes the outcome to True. \\
remove\_role\_revokes\_permission & Removing a role changes the outcome to False. \\
add\_permission & Adding a permission to a role grants access. \\
remove\_permission & Removing a permission revokes access. \\
invalid\_user\_returns\_false & Unknown user is denied, not an error. \\
user\_with\_no\_roles\_is\_denied & Empty role list means no access. \\
duplicate\_role\_assignment & Idempotency: assigning twice keeps one entry. \\
duplicate\_permission\_assignment & Idempotency: adding twice keeps one entry. \\
create\_user\_with\_unknown\_role & Validation: referencing a missing role raises. \\
assign\_unknown\_role & Validation: assigning a missing role raises. \\
api\_example\_scenario & End-to-end over HTTP for the assignment scenario. \\
api\_create\_user\_returns\_201 & Correct created status code. \\
api\_unknown\_role\_returns\_404 & Service error maps to HTTP 404. \\
api\_missing\_name\_returns\_422 & Pydantic validation maps to HTTP 422. \\
\end{longtable}

\section{Common Interview Questions About Testing}
\begin{itemize}[leftmargin=1.4em]
  \item \emph{Unit vs. integration?} Service tests are unit; TestClient tests are
        integration (they exercise routing, validation, serialisation).
  \item \emph{Why fixtures?} Fresh state per test ensures isolation and
        repeatability.
  \item \emph{How to test errors?} \texttt{pytest.raises(NotFoundError)} for the
        service; assert status codes for the API.
  \item \emph{What is not tested and why?} The frontend (out of scope for unit
        tests) and persistence durability (there is none by design).
\end{itemize}

% =========================================================================
\chapter{Scalability Discussion}
% =========================================================================

\emph{``How would you support 10 million users?''} The following are directions,
not implementations.

\begin{itemize}[leftmargin=1.4em]
  \item \textbf{Database.} Move from dicts to a real database. Tables:
        \texttt{users}, \texttt{roles}, \texttt{permissions}, and join tables
        \texttt{user\_roles}, \texttt{role\_permissions}. This gives durability,
        transactions, and shared state across processes. The first step (SQLite)
        is already implemented in \texttt{sql\_version/}; PostgreSQL is the
        production target.
  \item \textbf{Storage layer.} \texttt{sql\_version/} implements this with the
        standard library's \texttt{sqlite3} (plain, parameterized SQL) behind the
        same service interface, so the service code barely changes. An ORM such
        as SQLAlchemy is an optional later convenience if the SQL grows.
  \item \textbf{Indexes.} Index join-table foreign keys and the
        \texttt{(resource, action)} columns so permission queries stay fast.
  \item \textbf{Redis cache.} Cache each user's effective permission set;
        permission checks become O(1) cache hits. Invalidate on role/permission
        changes.
  \item \textbf{JWT + Authentication.} Users log in and receive a signed token
        carrying their id (and possibly cached roles). The service verifies the
        token instead of trusting a raw \texttt{user\_id}.
  \item \textbf{Authorization middleware.} A dependency/middleware that enforces
        required permissions per route.
  \item \textbf{Role hierarchy \& inheritance.} ``admin'' inherits ``editor''
        inherits ``viewer''. Compute the transitive closure of roles.
  \item \textbf{Permission groups \& wildcards.} Group related permissions;
        support \texttt{documents:*} to match any action.
  \item \textbf{Audit logs.} Append-only record of who was granted/checked what,
        for compliance and debugging.
  \item \textbf{Horizontal scaling.} With external storage the API is stateless,
        so run many instances behind a load balancer.
  \item \textbf{Microservices.} Extract an ``authz service'' other services call.
  \item \textbf{Event-driven architecture.} Publish role/permission-change events
        to invalidate caches and update read models asynchronously.
\end{itemize}

Because the service layer is isolated from HTTP and storage, most of these are
additive changes rather than rewrites.

% =========================================================================
\chapter{Interview Questions}
% =========================================================================

Roughly one hundred questions grouped by topic. Each has an ideal answer, a
common mistake, and a follow-up.

\section{RBAC and Design}

\textbf{Q1. What is RBAC?}
\qa{An authorization model where permissions are grouped into roles and roles are
assigned to users; access is decided by the roles a user holds.}
{Confusing authentication (who you are) with authorization (what you may do).}
{How does RBAC differ from ABAC?}

\textbf{Q2. RBAC vs ABAC vs ACL?}
\qa{ACL lists permissions per object; RBAC groups them into roles; ABAC decides
from attributes (user, resource, environment) via policies.}
{Claiming one is universally best.}
{When would you pick ABAC over RBAC?}

\textbf{Q3. What is a permission here?}
\qa{A (resource, action) pair, e.g. documents:write.}
{Adding an id to permissions when the pair already identifies them.}
{How would you support wildcard actions?}

\textbf{Q4. Why store role ids on users instead of role objects?}
\qa{So each role is a single source of truth; editing it affects all holders and
avoids duplicated, drifting copies.}
{Embedding full role objects and having to update every user on change.}
{What are the downsides of references?}

\textbf{Q5. Why are permissions stored in a set?}
\qa{O(1) membership and automatic de-duplication, which makes adding a duplicate a
no-op.}
{Using a list and getting duplicates or O(n) checks.}
{What must a type satisfy to live in a set?}

\textbf{Q6. Why is Permission frozen?}
\qa{Immutability makes it hashable, so it can be a set element and dict key.}
{Forgetting frozen and getting ``unhashable type''.}
{What does frozen generate under the hood?}

\textbf{Q7. What does default deny mean?}
\qa{If no role grants the permission, the answer is False; access is never granted
implicitly.}
{Returning True on error or for unknown users.}
{Where else should default deny apply?}

\textbf{Q8. How do you handle an unknown user in a check?}
\qa{Return False---an unknown subject can do nothing---matching the spec.}
{Raising 500 or leaking whether the user exists.}
{When might you prefer a 404 instead?}

\textbf{Q9. Is removing a role idempotent? Why does that matter?}
\qa{Yes; \texttt{set.discard} does nothing if absent. Idempotent DELETE is
RESTful and safe to retry.}
{Raising on removing a non-existent link.}
{Which HTTP methods should be idempotent?}

\textbf{Q10. How would you add role hierarchy?}
\qa{Give roles parent roles and compute the transitive closure of permissions when
checking (or precompute it).}
{Infinite loops on cyclic hierarchies.}
{How do you detect cycles?}

\textbf{Q11. How would wildcard permissions work?}
\qa{Match documents:* to any action on documents; check exact match first, then
wildcards.}
{Letting wildcards silently over-grant.}
{How do you keep wildcard checks fast?}

\textbf{Q12. Multi-tenant RBAC?}
\qa{Scope roles/permissions to a tenant id so identical role names don't collide
across organisations.}
{Global roles leaking across tenants.}
{Where does the tenant id come from at request time?}

\textbf{Q13. Where would you enforce least privilege?}
\qa{By designing narrow roles and granting only what a job needs; review
periodically.}
{One ``super'' role everyone gets.}
{How do you audit over-privileged users?}

\textbf{Q14. How do you answer ``who can delete invoices''?}
\qa{Find roles containing invoices:delete, then users holding those roles;
optionally index permission -> roles.}
{Scanning all users linearly every time.}
{What index makes this O(1)?}

\textbf{Q15. What are the limits of RBAC?}
\qa{Role explosion when access depends on many attributes; ABAC or hybrid models
handle that better.}
{Trying to encode every attribute as a role.}
{How do you mitigate role explosion?}

\section{Python}

\textbf{Q16. What is a dataclass?}
\qa{A decorator that auto-generates \texttt{\_\_init\_\_}, \texttt{\_\_repr\_\_},
\texttt{\_\_eq\_\_} from annotated fields.}
{Writing all boilerplate by hand.}
{What does \texttt{frozen=True} add?}

\textbf{Q17. Why default\_factory instead of a default list?}
\qa{A mutable default (\texttt{[]}) is shared across instances; \texttt{field(default\_factory=list)}
creates a fresh one each time.}
{Using a shared mutable default and getting cross-instance bugs.}
{Show the classic mutable-default trap.}

\textbf{Q18. set vs list vs dict complexity?}
\qa{Set/dict membership and dict lookup are O(1) average; list membership is
O(n).}
{Assuming list \texttt{in} is fast.}
{When is a list still the right choice?}

\textbf{Q19. What makes an object hashable?}
\qa{Stable \texttt{\_\_hash\_\_} and \texttt{\_\_eq\_\_}; immutable value objects
qualify.}
{Putting a mutable object in a set.}
{Why must equal objects have equal hashes?}

\textbf{Q20. What are type hints and do they run?}
\qa{Annotations for readability and tooling; not enforced at runtime by Python
itself, but used by Pydantic/FastAPI.}
{Believing hints enforce types at runtime.}
{How does Pydantic use them?}

\textbf{Q21. dict.get vs indexing?}
\qa{\texttt{get} returns \texttt{None} (or a default) if absent; \texttt{[]}
raises \texttt{KeyError}.}
{Using \texttt{[]} where absence is expected.}
{When do you want the KeyError?}

\textbf{Q22. What is a generator/comprehension?}
\qa{Concise iterable construction; comprehensions build lists/sets, generators
yield lazily.}
{Building a huge list when a generator suffices.}
{When does laziness matter here?}

\textbf{Q23. What is \texttt{is} vs \texttt{==}?}
\qa{\texttt{is} tests identity; \texttt{==} tests value equality. We use
\texttt{is None} for identity.}
{Using \texttt{== None}.}
{Why is \texttt{is None} preferred?}

\textbf{Q24. Truthiness of empty containers?}
\qa{Empty set/list/dict/string are falsy.}
{Comparing \texttt{len(x) == 0} unnecessarily.}
{Any pitfalls with \texttt{0} vs empty?}

\textbf{Q25. What is a virtual environment?}
\qa{An isolated per-project set of dependencies so projects don't clash.}
{Installing globally and getting version conflicts.}
{How does \texttt{requirements.txt} fit in?}

\textbf{Q26. How do custom exceptions help?}
\qa{They carry structured context (entity, id) and can be handled specifically,
e.g. mapped to 404.}
{Raising bare \texttt{Exception} everywhere.}
{Where is \texttt{NotFoundError} translated?}

\textbf{Q27. What is \texttt{uuid4} and why use it?}
\qa{A random 128-bit id; collision-resistant without a central counter.}
{Using incremental ints that leak volume/order.}
{UUID vs auto-increment trade-offs?}

\textbf{Q28. GIL---does it matter here?}
\qa{The GIL serialises Python bytecode; our work is trivial and I/O-bound under
the server, so it's not a bottleneck at this scale.}
{Claiming threads give CPU parallelism in CPython.}
{How would you scale CPU-bound work?}

\textbf{Q29. Shallow vs deep copy---relevant?}
\qa{We don't copy domain objects; users reference roles by id, so there's no
aliasing bug.}
{Copying role objects into users.}
{When would a copy be needed?}

\textbf{Q30. What is \texttt{\_\_init\_\_.py} for?}
\qa{Marks a directory as a package so \texttt{import app...} works.}
{Forgetting it and hitting import errors.}
{How do namespace packages differ?}

\section{FastAPI}

\textbf{Q31. How does FastAPI validate input?}
\qa{It parses the JSON body into the Pydantic model in the signature; failures
become a 422 automatically.}
{Manually validating inside the handler.}
{What does the 422 body contain?}

\textbf{Q32. What is a response\_model?}
\qa{The declared output schema; it serialises and filters the return value and
documents the endpoint.}
{Returning internal objects with extra fields.}
{Can it strip secrets from output?}

\textbf{Q33. What is Depends?}
\qa{Dependency injection: FastAPI calls the dependency and passes its result;
enables shared setup and test overrides.}
{Using a bare global that can't be swapped in tests.}
{How do you override a dependency in tests?}

\textbf{Q34. APIRouter vs app?}
\qa{Router groups endpoints in a module; \texttt{include\_router} mounts them,
keeping \texttt{main.py} thin.}
{Defining everything on the app object.}
{Can routers have prefixes/tags?}

\textbf{Q35. How are path vs body params distinguished?}
\qa{Params named in the path are path params; a Pydantic-typed param is the body;
scalars become query params.}
{Expecting a body on a GET.}
{How to send a body on DELETE?}

\textbf{Q36. Sync vs async endpoints?}
\qa{FastAPI supports both; sync handlers run in a threadpool. Our logic is fast
and non-blocking, so sync is fine.}
{Doing blocking I/O in an async handler.}
{When must a handler be async?}

\textbf{Q37. Where do the docs come from?}
\qa{FastAPI generates an OpenAPI schema from routes and models, served at
\texttt{/docs} and \texttt{/redoc}.}
{Writing API docs by hand.}
{How do examples get into the docs?}

\textbf{Q38. What is CORS and why enable it?}
\qa{Browsers block cross-origin requests unless the server opts in; we enable it
so a separately served frontend can call the API.}
{Debugging a ``CORS error'' as a backend bug.}
{Why is allow-all unsafe in production?}

\textbf{Q39. How do you return a custom status code?}
\qa{\texttt{status\_code=status.HTTP\_201\_CREATED} on the decorator, or raise an
exception mapped to a code.}
{Returning 200 for a creation.}
{201 vs 200 vs 204?}

\textbf{Q40. What is an exception handler?}
\qa{\texttt{@app.exception\_handler(NotFoundError)} centralises turning an
exception into an HTTP response.}
{Repeating try/except in every route.}
{Why is central handling cleaner?}

\textbf{Q41. How does FastAPI serialise a set?}
\qa{We convert sets to sorted lists in the response helpers so JSON is valid and
deterministic.}
{Trying to JSON-encode a raw set.}
{Why sort the output?}

\textbf{Q42. What runs the app?}
\qa{Uvicorn, an ASGI server; \texttt{uvicorn app.main:app} imports the app
object.}
{Confusing WSGI (sync) with ASGI (async).}
{What does \texttt{--reload} do?}

\textbf{Q43. What is StaticFiles used for here?}
\qa{Serving the frontend from the same process at \texttt{/}, so the demo needs
one server.}
{Route ordering hiding API paths behind the static mount.}
{Why mount static last?}

\textbf{Q44. How do you version an API?}
\qa{Prefix routes (e.g. \texttt{/v1}) or use headers; keep old versions until
clients migrate.}
{Breaking clients with silent changes.}
{Where would you add the prefix here?}

\textbf{Q45. What is dependency\_overrides?}
\qa{A dict on the app mapping a dependency to a replacement, used to inject a
clean service in tests.}
{Sharing seeded state across tests.}
{Why clear it after each test?}

\section{REST APIs}

\textbf{Q46. What makes an API RESTful?}
\qa{Resources addressed by URLs, standard verbs, statelessness, and appropriate
status codes.}
{Using GET to mutate state.}
{Is this API fully REST or pragmatic?}

\textbf{Q47. Which verb for create/read/update/delete?}
\qa{POST/GET/PUT-PATCH/DELETE respectively.}
{Using POST for everything.}
{PUT vs PATCH?}

\textbf{Q48. Why 201 vs 200?}
\qa{201 signals a new resource was created; 200 a successful general response.}
{Returning 200 on create.}
{What should a 201 body contain?}

\textbf{Q49. Why 404 vs 422?}
\qa{404: referenced resource missing; 422: the request itself is malformed.}
{Returning 400/500 for validation errors.}
{Where does each occur in this project?}

\textbf{Q50. Are your endpoints idempotent?}
\qa{GET/DELETE and the set-based POSTs are idempotent; creating a new user is
not.}
{Assuming all POSTs are idempotent.}
{How would you make create idempotent?}

\textbf{Q51. Statelessness---do you comply?}
\qa{Each request carries all needed data; the server keeps no per-client session.
State lives in the shared store.}
{Confusing app state with session state.}
{How does statelessness help scaling?}

\textbf{Q52. How would you paginate GET /users?}
\qa{Add \texttt{limit}/\texttt{offset} or cursor query params.}
{Returning millions of rows at once.}
{Offset vs cursor pagination?}

\textbf{Q53. How to handle partial failure on bulk assign?}
\qa{Validate all first (all-or-nothing) or return per-item results; be explicit.}
{Half-applying and leaving inconsistent state.}
{Which does \texttt{create\_user} choose?}

\textbf{Q54. Content negotiation?}
\qa{We use JSON only; \texttt{Content-Type: application/json} in and out.}
{Ignoring the header entirely.}
{How would you add CSV export?}

\textbf{Q55. Why send the permission in a DELETE body?}
\qa{A permission has no id, so the (resource, action) pair identifies it; DELETE
with a body is allowed.}
{Assuming DELETE cannot carry a body.}
{Alternative: encode it in the path?}

\section{Testing}

\textbf{Q56. What does a good unit test isolate?}
\qa{One unit of behaviour with fresh state and no external dependencies.}
{Tests that depend on each other's order.}
{How do fixtures ensure isolation?}

\textbf{Q57. What is a fixture?}
\qa{A reusable setup function (e.g. a fresh service) injected into tests.}
{Copy-pasting setup into every test.}
{What are fixture scopes?}

\textbf{Q58. How do you test that an error is raised?}
\qa{\texttt{with pytest.raises(NotFoundError): ...}}
{Wrapping in try/except and forgetting to fail if no error.}
{How to assert on the message?}

\textbf{Q59. Unit vs integration here?}
\qa{Service tests are unit; TestClient tests integrate routing, validation, and
serialisation.}
{Calling everything ``unit''.}
{Which catches a wrong status code?}

\textbf{Q60. What is TestClient?}
\qa{An in-process HTTP client (httpx-based) that drives the ASGI app without a
real server.}
{Spinning up Uvicorn for tests.}
{Why is it faster than a live server?}

\textbf{Q61. What is coverage and its limit?}
\qa{The percent of code executed by tests; high coverage doesn't guarantee correct
assertions.}
{Chasing 100\% with assertion-free tests.}
{What would you prioritise covering?}

\textbf{Q62. How would you test idempotency?}
\qa{Perform the action twice and assert the state equals the single-action state.}
{Only running the action once.}
{Which tests do this here?}

\textbf{Q63. Parametrized tests?}
\qa{\texttt{@pytest.mark.parametrize} runs one test over many inputs.}
{Copying a test many times.}
{Where would it fit here?}

\textbf{Q64. What is a regression test?}
\qa{A test added to lock in a bug fix so it can't return.}
{Fixing a bug without a test.}
{Give an example for this service.}

\textbf{Q65. Mock vs real object?}
\qa{Mock external/slow dependencies; here the store is in-memory and fast, so we
use the real thing.}
{Over-mocking and testing mocks, not code.}
{When would you mock storage?}

\section{Security}

\textbf{Q66. AuthN vs AuthZ?}
\qa{Authentication verifies identity; authorization (this service) decides
permitted actions.}
{Using the terms interchangeably.}
{Which does this project implement?}

\textbf{Q67. Why is trusting a raw user\_id insecure?}
\qa{A caller could impersonate anyone; identity must come from a verified token,
not the request body.}
{Shipping this as production auth.}
{How would JWT fix it?}

\textbf{Q68. What is a JWT?}
\qa{A signed token carrying claims (e.g. user id); the server verifies the
signature without a session store.}
{Storing secrets in the payload (it's readable).}
{Access vs refresh tokens?}

\textbf{Q69. Why default deny is a security property?}
\qa{Failing closed means bugs tend to deny rather than over-grant.}
{Failing open on errors.}
{Where do we fail closed?}

\textbf{Q70. Risks of \texttt{allow\_origins=["*"]}?}
\qa{Any site can call the API from a browser; fine for a local demo, unsafe with
credentials in production.}
{Leaving it open in production.}
{How do you scope origins?}

\textbf{Q71. What is privilege escalation?}
\qa{Gaining access beyond one's role; RBAC changes must be themselves
permission-protected.}
{Letting any user assign roles.}
{Who should call the role endpoints?}

\textbf{Q72. Why audit logs?}
\qa{They record who changed or checked access, essential for compliance and
incident response.}
{No trail after a breach.}
{What fields would you log?}

\textbf{Q73. Input validation as defense?}
\qa{Pydantic rejects malformed input at the edge, shrinking the attack surface.}
{Trusting client input.}
{What can't validation catch?}

\textbf{Q74. Sensitive data in responses?}
\qa{\texttt{response\_model} whitelists fields, preventing accidental leaks.}
{Returning raw internal objects.}
{Example of a field to hide?}

\textbf{Q75. Rate limiting---why?}
\qa{Prevents abuse and brute force; apply per client/IP.}
{No limits, enabling credential stuffing.}
{Where would you enforce it?}

\section{System Design and Scalability}

\textbf{Q76. First step to scale past one process?}
\qa{Externalise state into a shared database so multiple stateless instances
agree.}
{Scaling instances while state stays in-memory.}
{What breaks with in-memory at scale?}

\textbf{Q77. Schema for RBAC in SQL?}
\qa{users, roles, permissions, and join tables user\_roles, role\_permissions,
with indexes on foreign keys.}
{One giant denormalised table.}
{How do joins answer a permission check?}

\textbf{Q78. How does caching speed checks?}
\qa{Precompute a user's effective permission set into Redis; a check is one O(1)
lookup.}
{Never invalidating the cache.}
{When do you invalidate it?}

\textbf{Q79. Cache invalidation strategy?}
\qa{Invalidate a user's cached set when their roles change, and all holders' sets
when a role's permissions change.}
{Forgetting the role-change fan-out.}
{TTL vs explicit invalidation?}

\textbf{Q80. How to make the API horizontally scalable?}
\qa{Keep it stateless with external storage, then run N instances behind a load
balancer.}
{Sticky in-memory sessions.}
{What about cache consistency across instances?}

\textbf{Q81. When to split into microservices?}
\qa{When authz has independent scaling/ownership needs; expose an internal authz
service.}
{Premature microservices adding latency.}
{Trade-offs of the split?}

\textbf{Q82. Role hierarchy at scale?}
\qa{Precompute effective permissions via the transitive closure and cache it;
recompute on hierarchy edits.}
{Recomputing closure on every request.}
{How to detect cycles cheaply?}

\textbf{Q83. Event-driven invalidation?}
\qa{Publish role/permission-change events; consumers refresh caches and read
models.}
{Synchronous fan-out blocking writes.}
{At-least-once vs exactly-once?}

\textbf{Q84. Read vs write ratio implications?}
\qa{Checks (reads) vastly outnumber writes, so optimise reads with caching and
denormalised permission sets.}
{Optimising the rare write path.}
{What consistency can you relax?}

\textbf{Q85. How to answer ``who can do X'' at scale?}
\qa{Maintain a reverse index permission -> roles -> users, updated on writes.}
{Full scans per query.}
{Where is that index stored?}

\textbf{Q86. Indexing basics?}
\qa{Indexes turn O(n) scans into O(log n) lookups on indexed columns at some write
cost.}
{Indexing everything indiscriminately.}
{Which columns here?}

\textbf{Q87. Handling hot keys (a huge role)?}
\qa{A role held by millions makes invalidation expensive; batch or lazily
recompute, and cache with TTL.}
{Synchronously updating millions on one edit.}
{How to bound the fan-out?}

\textbf{Q88. Backwards-compatible migrations?}
\qa{Additive changes, dual-write/read during transition, then remove old paths.}
{Big-bang breaking migrations.}
{How to roll back safely?}

\textbf{Q89. Observability?}
\qa{Structured logs, metrics (check latency, deny rate), and tracing across
layers.}
{No visibility into authz decisions.}
{What metric signals a problem?}

\textbf{Q90. Where is the bottleneck likely?}
\qa{The permission check on the read path; solve with caching and precomputation.}
{Assuming the database write path is hottest.}
{How do you confirm with data?}

\section{Mixed / Rapid Fire}

\textbf{Q91. Why separate schemas from models?}
\qa{So the wire format can evolve independently of internal storage and to avoid
leaking internals.}
{Reusing one class for both and coupling them.}
{Give a case where they diverge.}

\textbf{Q92. Why a service layer at all?}
\qa{It centralises rules, stays framework-free, and is unit-testable without
HTTP.}
{Putting logic in route handlers.}
{What would break if logic lived in routes?}

\textbf{Q93. What is the single source of truth for a role?}
\qa{The \texttt{roles} dict entry; users only hold ids.}
{Copying role data into users.}
{Why does this simplify updates?}

\textbf{Q94. How is determinism achieved in responses?}
\qa{Helpers sort role ids and permissions before returning.}
{Relying on set iteration order.}
{Why does sorting help tests?}

\textbf{Q95. What happens on server restart?}
\qa{In-memory data resets to the seed; by design there is no persistence.}
{Expecting durability.}
{What's the first fix for that?}

\textbf{Q96. How would you add Delete User?}
\qa{A service method popping the user from the dict and a \texttt{DELETE /users/\{id\}}
route returning 204.}
{Leaving dangling references.}
{Any cleanup needed? (None: users hold the ids.)}

\textbf{Q97. How would you add logging?}
\qa{Configure the \texttt{logging} module or middleware to record requests and
authz decisions.}
{Using \texttt{print}.}
{What not to log (secrets)?}

\textbf{Q98. What is the time complexity of the whole check path?}
\qa{O(R) in the user's roles, dominated by set lookups that are O(1) each.}
{Claiming O(1) without noting the role loop.}
{How to reach O(1)?}

\textbf{Q99. How do you keep the project explainable?}
\qa{Clear layering, small files, descriptive names, and minimal abstraction.}
{Adding patterns that aren't needed.}
{Which file would you show first?}

\textbf{Q100. What would you build next given more time?}
\qa{Persistence + cache, authentication, and delete endpoints, in that order.}
{Jumping to microservices first.}
{Why that order?}

% =========================================================================
\chapter{Feature Extension Guide}
% =========================================================================

For each feature: high-level thinking, files to touch, functions that change,
pseudocode, and expected complexity. \emph{Discussion and pseudocode only---not
implemented in the project.}

\section{Delete User / Delete Role}
\emph{Note: this extension is already implemented in the project (see
\texttt{delete\_user}/\texttt{delete\_role} in \texttt{services.py} and the
DELETE routes), so it doubles as a worked example of the pattern below.}\\
\textbf{Thinking.} Symmetric to create. Deleting a role should also drop it from
every user's \texttt{role\_ids}.\\
\textbf{Files.} \texttt{services.py}, \texttt{routes.py}.\\
\textbf{Functions.} add \texttt{delete\_user}, \texttt{delete\_role}; new DELETE
routes.
\begin{lstlisting}[language=Python]
def delete_role(self, role_id):
    self.get_role(role_id)                 # 404 if missing
    del self.store.roles[role_id]
    for user in self.store.users.values(): # cascade
        user.role_ids.discard(role_id)
\end{lstlisting}
\textbf{Complexity.} \texttt{delete\_user} O(1); \texttt{delete\_role} O(U).

\section{Authentication / JWT / OAuth}
\textbf{Thinking.} Identity must be verified, not trusted from the body. Issue a
signed token on login; verify it on each request; derive \texttt{user\_id} from
the token.\\
\textbf{Files.} new \texttt{auth.py}, \texttt{main.py}, \texttt{routes.py}.\\
\textbf{Functions.} \texttt{create\_token}, \texttt{get\_current\_user}
dependency; protect routes with it.
\begin{lstlisting}[language=Python]
def get_current_user(token = Depends(oauth2_scheme)):
    payload = decode_and_verify(token)     # raises 401 if invalid
    return payload["sub"]                  # the user id
\end{lstlisting}
\textbf{Complexity.} O(1) signature verification per request.

\section{Admin Dashboard}
\textbf{Thinking.} A richer frontend (still vanilla JS or a framework) calling the
existing endpoints plus new delete/edit routes.\\
\textbf{Files.} \texttt{frontend/}, possibly new routes.\\
\textbf{Functions.} none in the core; mostly UI.\\
\textbf{Complexity.} UI-bound.

\section{SQL / NoSQL Persistence (behind the same interface)}
\textbf{Thinking.} Implement a store with the same methods as \texttt{InMemoryStore}
backed by a database; the service is unchanged.\\
\textbf{Files.} new \texttt{sql\_store.py}, wire it in \texttt{dependencies.py}.\\
\textbf{Functions.} same surface (\texttt{users}, \texttt{roles} access) via
repository methods.
\begin{lstlisting}[language=Python]
class SqlStore:
    def get_user(self, user_id): ...   # SELECT ... WHERE id = ?
    def save_user(self, user): ...     # INSERT/UPDATE
\end{lstlisting}
\textbf{Complexity.} O(log n) indexed lookups; network latency per call.

\section{Caching}
\textbf{Thinking.} Cache each user's effective permission set; check hits the
cache first.\\
\textbf{Files.} \texttt{services.py} (read-through + invalidation).\\
\textbf{Functions.} \texttt{can\_user\_perform\_action} (read cache);
\texttt{assign\_role}/\texttt{remove\_role}/\texttt{add\_permission}/\texttt{remove\_permission}
(invalidate).
\begin{lstlisting}[language=Python]
def can_user_perform_action(self, user_id, resource, action):
    perms = cache.get(user_id) or self._recompute_and_cache(user_id)
    return Permission(resource, action) in perms
\end{lstlisting}
\textbf{Complexity.} O(1) on cache hit; O(R*P) to recompute on miss.

\section{Overwrite / List of Permissions}
\textbf{Thinking.} Add a \texttt{set\_permissions} to replace a role's set, and
accept a list of permissions in add/remove.\\
\textbf{Files.} \texttt{schemas.py}, \texttt{services.py}, \texttt{routes.py}.\\
\textbf{Functions.} \texttt{set\_permissions(role\_id, permissions)}.
\begin{lstlisting}[language=Python]
def set_permissions(self, role_id, permissions):
    role = self.get_role(role_id)
    role.permissions = set(permissions)    # full overwrite
    return role
\end{lstlisting}
\textbf{Complexity.} O(P).

\section{Role Hierarchy / Permission Inheritance}
\textbf{Thinking.} Roles gain \texttt{parent\_ids}; effective permissions are the
union over the role and its ancestors.\\
\textbf{Files.} \texttt{models.py}, \texttt{services.py}.\\
\textbf{Functions.} \texttt{effective\_permissions(role)} with cycle detection.
\begin{lstlisting}[language=Python]
def effective_permissions(self, role_id, seen=None):
    seen = seen or set()
    if role_id in seen: return set()       # cycle guard
    seen.add(role_id)
    role = self.get_role(role_id)
    perms = set(role.permissions)
    for parent in role.parent_ids:
        perms |= self.effective_permissions(parent, seen)
    return perms
\end{lstlisting}
\textbf{Complexity.} O(V+E) over the role graph; cache the result.

\section{Permission Groups}
\textbf{Thinking.} A named group expands to a set of permissions; roles reference
groups.\\
\textbf{Files.} \texttt{models.py}, \texttt{services.py}.\\
\textbf{Functions.} expand groups when computing effective permissions.\\
\textbf{Complexity.} O(size of expansion).

\section{Wildcard Permissions}
\textbf{Thinking.} Support \texttt{resource:*}. On check, test the exact
permission and any wildcard forms.\\
\textbf{Files.} \texttt{services.py}.\\
\textbf{Functions.} \texttt{can\_user\_perform\_action}.
\begin{lstlisting}[language=Python]
candidates = {Permission(resource, action), Permission(resource, "*")}
return any(c in perms for c in candidates)
\end{lstlisting}
\textbf{Complexity.} O(1) extra lookups.

\section{Multi-Tenant RBAC}
\textbf{Thinking.} Add \texttt{tenant\_id} to users/roles; scope all lookups and
checks by tenant.\\
\textbf{Files.} \texttt{models.py}, \texttt{storage.py}, \texttt{services.py}.\\
\textbf{Functions.} every query gains a tenant filter.\\
\textbf{Complexity.} unchanged; keys become (tenant, id).

\section{Logging}
\textbf{Thinking.} Standard \texttt{logging} plus middleware for request/decision
logs.\\
\textbf{Files.} \texttt{main.py}.\\
\textbf{Complexity.} negligible.

\section{Rate Limiting}
\textbf{Thinking.} Middleware counting requests per client in a window (token
bucket), often backed by Redis.\\
\textbf{Files.} \texttt{main.py} (middleware).\\
\textbf{Complexity.} O(1) per request.

\section{Docker / CI-CD / Kubernetes}
\textbf{Thinking.} A \texttt{Dockerfile} runs Uvicorn; CI runs \texttt{pytest} on
push; Kubernetes runs multiple stateless replicas with a shared store.\\
\textbf{Files.} \texttt{Dockerfile}, CI config, k8s manifests.\\
\textbf{Complexity.} operational, not algorithmic.

% =========================================================================
\chapter{Code Reading Guide}
% =========================================================================

\section{Recommended Reading Order}
\begin{enumerate}[leftmargin=1.6em]
  \item \texttt{app/models.py} --- the vocabulary (User, Role, Permission).
  \item \texttt{app/storage.py} --- where they live (two dicts).
  \item \texttt{app/services.py} --- the rules and the permission check.
  \item \texttt{app/schemas.py} --- the API's input/output shapes.
  \item \texttt{app/routes.py} --- how HTTP maps to service calls.
  \item \texttt{app/dependencies.py} and \texttt{app/main.py} --- wiring and app.
  \item \texttt{app/seed.py} --- the example data.
  \item \texttt{tests/test\_rbac.py} --- behaviour, as executable documentation.
  \item \texttt{frontend/} --- the demo UI.
\end{enumerate}

\section{How to Explain the Project in an Interview}
\begin{enumerate}[leftmargin=1.6em]
  \item \textbf{One sentence:} ``An RBAC service---users have roles, roles have
        permissions, and one function answers can-user-perform-action.''
  \item \textbf{Draw the three boxes} (User, Role, Permission) and the layers
        (routes, service, storage).
  \item \textbf{Walk one check} end to end (frontend to storage and back).
  \item \textbf{Justify two choices:} references over embedding, and sets over
        lists.
  \item \textbf{State the complexity} (O(R)) and how you'd cache to O(1).
  \item \textbf{Pivot to scale} using the scalability chapter.
\end{enumerate}

\section{What Interviewers Are Evaluating}
\begin{itemize}[leftmargin=1.4em]
  \item \textbf{Correctness:} does the check match the spec and the examples?
  \item \textbf{Code quality:} clear names, small functions, separation of
        concerns.
  \item \textbf{Design:} sensible data structures and layering.
  \item \textbf{Completeness:} tests, docs, run instructions, seed data.
  \item \textbf{Scalability thinking:} can you articulate the path to millions of
        users without over-building now?
\end{itemize}

% =========================================================================
\chapter{Cheat Sheet}
% =========================================================================

\section{FastAPI}
\begin{lstlisting}[language=Python]
from fastapi import FastAPI, APIRouter, Depends, status
app = FastAPI()
router = APIRouter()

@router.post("/users", response_model=UserResponse,
             status_code=status.HTTP_201_CREATED)
def create_user(body: CreateUserRequest,
                service: RBACService = Depends(get_service)):
    return _user_to_response(service.create_user(body.name, body.role_ids))

app.include_router(router)

@app.exception_handler(NotFoundError)
def handle(request, exc):
    return JSONResponse(status_code=404, content={"detail": str(exc)})
\end{lstlisting}

\section{Pydantic v2}
\begin{lstlisting}[language=Python]
from pydantic import BaseModel, Field

class CreateUserRequest(BaseModel):
    name: str = Field(min_length=1)
    role_ids: list[str] = Field(default_factory=list)

# Parse + validate:  CreateUserRequest(**payload)
# Serialize:         model.model_dump()
\end{lstlisting}

\section{Python Classes / Dataclasses}
\begin{lstlisting}[language=Python]
from dataclasses import dataclass, field

@dataclass(frozen=True)   # immutable -> hashable -> set/dict key
class Permission:
    resource: str
    action: str

@dataclass
class Role:
    id: str
    name: str
    permissions: set[Permission] = field(default_factory=set)
\end{lstlisting}

\section{REST Quick Reference}
\begin{longtable}{p{0.16\textwidth} p{0.24\textwidth} p{0.5\textwidth}}
\textbf{Verb} & \textbf{Use} & \textbf{Idempotent?} \\
\hline
\endhead
GET & read & yes \\
POST & create / action & no (create) \\
PUT & replace & yes \\
PATCH & partial update & no (typically) \\
DELETE & remove & yes \\
\end{longtable}
Status codes: 200 OK, 201 Created, 204 No Content, 400 Bad Request, 401
Unauthorized, 403 Forbidden, 404 Not Found, 422 Unprocessable Entity, 500 Server
Error.

\section{RBAC Concepts}
\begin{itemize}[leftmargin=1.4em]
  \item User \texttt{->} Roles \texttt{->} Permissions (resource, action).
  \item Default deny; grant only via a matching permission.
  \item References (ids), not embedded copies, for a single source of truth.
  \item Idempotent grants/removals via sets.
  \item AuthN (identity) is separate from AuthZ (this service).
\end{itemize}

\section{Time Complexities}
\begin{longtable}{p{0.55\textwidth} p{0.4\textwidth}}
\textbf{Operation} & \textbf{Cost} \\
\hline
\endhead
dict lookup / set membership (avg) & O(1) \\
create role with P permissions & O(P) \\
create user with R roles (validate) & O(R) \\
assign/remove role, add/remove permission & O(1) \\
\texttt{can\_user\_perform\_action} & O(R) \\
delete role (cascade to users) & O(U) \\
effective permissions with hierarchy & O(V+E) \\
\end{longtable}

\section{Common Interview One-Liners}
\begin{itemize}[leftmargin=1.4em]
  \item ``Permissions are a (resource, action) pair; roles bundle them; users
        hold roles.''
  \item ``Users store role ids so a role is a single source of truth.''
  \item ``Sets give O(1) checks and idempotent, duplicate-free assignments.''
  \item ``The check is O(R); cache the effective set to make it O(1).''
  \item ``Layers---routes, service, storage---so I can swap storage or add auth
        without touching the core logic.''
  \item ``Unknown user returns False: fail closed, default deny.''
\end{itemize}

\vspace{1em}
\begin{center}
\textit{End of report.}
\end{center}

\end{document}
